\chapter{Troubleshooting}\label{ch:troubleshooting}

This chapter is organized by what you see on the screen. Find the situation
you are in, then look for the message or behavior that matches yours. Each
entry says why it happens, what to do, and where the full explanation is.
Limitations that have no fix yet are collected in
Chapter~\ref{ch:knownissues}.

\section{First things to check}\label{sec:troubleshooting-first}

Many problems have one of a few simple causes. Check these first.

\begin{enumerate}
  \item \textbf{Is the right environment active?} With the conda installer
  (Option~B), type \texttt{conda activate ProPrep} before starting ProPrep.
  Then type \texttt{echo \$AMBERHOME}. The answer should be empty or a folder
  inside \file{ProPrep}. Any other folder means ProPrep will use another Amber
  installation (Chapter~\ref{ch:installation}).
  \item \textbf{Are you online?} Downloading structures, AlphaFold models,
  BLAST searches and the OPM membrane check all need the internet. Some
  tools report a lost connection with a misleading message.
  \item \textbf{Are you in the right project folder?} ProPrep looks for
  session files only in the folder you start it from. Inside ProPrep, read
  the \prompt{Using ... file:} line at the top of each tool: it tells you
  which file the tool is working on (Chapter~\ref{ch:projectfolder}).
  \item \textbf{Look at the log file.} When a topology build fails, the real
  reason is often only in \file{leap.log} in the project folder. A failed
  batch run writes its own \file{error.log}. An MD stage writes
  \file{simulation.stderr} and an \file{.mdout} file in its stage folder.
  \item \textbf{Gave a wrong answer?} Type \typed{undo} at the next question
  to go back (Chapter~\ref{ch:sessions}).
  \item \textbf{ProPrep stopped or the window closed?} Your answers are saved
  as you go. Start again with \texttt{--resume-session} and the session
  file, and choose \textsf{Continue from end} (Chapter~\ref{ch:sessions}).
\end{enumerate}

\section{Installing, starting and moving around}\label{sec:troubleshooting-start}

{\bfseries Every start prints \prompt{MODELLER not found ... To install: conda
install -c salilab modeller}, although MODELLER is installed.}\\*
\emph{Why:} In version 1.23.0 this message also appears when only the MODELLER
license key is missing, which is the usual case.\\
\emph{What to do:} Do not install MODELLER again. Add your key to
\file{~/.proprep/modeller_key}. Without a key, everything works except loop,
residue and atom repair, standard mutations and homology models. See
Chapter~\ref{ch:installation}.

{\bfseries The installer stops with \prompt{No such file or directory},
\prompt{Permission denied}, \prompt{File or directory already exists},
\prompt{No space left on device} or \prompt{conda is not installed}.}\\*
\emph{Why:} A mistyped file name, a missing \texttt{bash} at the start of the
line, an existing installation, too little disk space, or (Option~B) no
conda.\\
\emph{What to do:} The table at the end of Chapter~\ref{ch:installation}
gives the fix for each message.

{\bfseries \texttt{proprep: command not found} (Option~B).}\\*
\emph{Why:} The ProPrep environment is not active in this Terminal window.\\
\emph{What to do:} Type \texttt{conda activate ProPrep}, then start ProPrep
again. See Chapter~\ref{ch:installation}.

{\bfseries tLEaP or sander messages show a folder from another installation,
for example \file{.../envs/etanalyze_env/bin/teLeap}.}\\*
\emph{Why:} ProPrep keeps an \texttt{AMBERHOME} setting that is already set in
the Terminal, so it uses that Amber instead of its own. It gives no warning.\\
\emph{What to do:} Exit ProPrep. Type \texttt{echo \$AMBERHOME}. If it names
another folder, type \texttt{unset AMBERHOME} (Option~A) or
\texttt{source \$CONDA\_PREFIX/amber.sh} after \texttt{conda activate ProPrep}
(Option~B), then start ProPrep from the same window. See
Chapter~\ref{ch:installation}.

{\bfseries The browser says \prompt{This site can't be reached} at
\texttt{127.0.0.1:8000}.}\\*
\emph{Why:} \texttt{proprep-web} is not running, or it chose another port
because 8000 was busy.\\
\emph{What to do:} Start \texttt{proprep-web} again and use the address it
prints, for example \prompt{(port 8420 was busy; using 8421)}. See
Chapter~\ref{ch:starting}.

{\bfseries Pressing \key{Enter} at \prompt{Select menu mode [1/2]} only repeats
the question.}\\*
\emph{Why:} This question has no default.\\
\emph{What to do:} Type \key{2} for the Full Menu used in this manual, or
start with \texttt{proprep --full-menu}. See Chapter~\ref{ch:starting}.

{\bfseries ProPrep stops at start-up with \prompt{Error: [Errno 2] No such
file or directory: 'sub/deeper'}.}\\*
\emph{Why:} The project name contains a \texttt{/}, which asks for two folder
levels; ProPrep can create only one.\\
\emph{What to do:} Use a simple project name made of letters, digits,
\texttt{-} and \texttt{\_}, without spaces. See Chapter~\ref{ch:starting}.

{\bfseries A project folder called \file{undo} appeared.}\\*
\emph{Why:} \typed{undo} typed at the project-name question is taken as the
name, because recording starts only after that question.\\
\emph{What to do:} Delete the folder and start again. \typed{undo} works from
the next question on. See Chapter~\ref{ch:sessions}.

{\bfseries ProPrep does not offer to resume your project (no
\prompt{Previous session files found!}).}\\*
\emph{Why:} ProPrep looks for session files only in the folder it was started
from, not in the project folders below it.\\
\emph{What to do:} Move into the project folder first (\texttt{cd
myproject}), then start ProPrep. See Chapter~\ref{ch:starting}.

{\bfseries You chose to resume, pressed \key{Enter} in the list of sessions,
and got a fresh session.}\\*
\emph{Why:} The default of \prompt{Select session to resume (0 to cancel)} is
\key{0}, which means cancel.\\
\emph{What to do:} Type the number of the session. See
Chapter~\ref{ch:sessions}.

{\bfseries ProPrep ended when you closed the browser tab.}\\*
\emph{Why:} \texttt{proprep-web} shuts itself down about two seconds after the
last tab closes.\\
\emph{What to do:} Nothing is lost: your files and the session file stay in
the project folder. Next time, exit with \key{x} first and download the
project before you close the tab. See Chapter~\ref{ch:starting}.

{\bfseries \prompt{Invalid choice: W. Please try again.} for a menu code
you know exists.}\\*
\emph{Why:} The Full Menu is case-sensitive.\\
\emph{What to do:} Type codes in lower case: \key{w}, \key{4e}. See
Chapter~\ref{ch:menus}.

{\bfseries \prompt{Received '1a', which is not one of the available options}
inside a tool.}\\*
\emph{Why:} Main-menu codes do not work inside a tool's own menu.\\
\emph{What to do:} Type \key{m} to go to the main menu first, then the code.
See Chapter~\ref{ch:menus}.

{\bfseries \prompt{Received '', which is not one of the available
options}.}\\*
\emph{Why:} You pressed \key{Enter} at a question that has no default.\\
\emph{What to do:} Type one of the letters or numbers offered. See
Chapter~\ref{ch:menus}.

{\bfseries A code from the Help page (\key{h}) opens the wrong tool, for
example \key{3a} opens the Homology Searcher.}\\*
\emph{Why:} In version 1.23.0 the Help page lists the menu codes of an older
version.\\
\emph{What to do:} Use the codes on the main menu itself, or the menu map in
Appendix~\ref{app:menumap}.

\section{Loading structures}\label{sec:troubleshooting-loading}

{\bfseries A tool works on an old or unexpected file, for example
\prompt{Using H-Stripped PDB file:} after you loaded a new structure.}\\*
\emph{Why:} Each tool picks its input from the project folder by a fixed order
of preference, not by date. Files ProPrep made (without hydrogens, assembly,
aligned, filtered, repaired, oriented) come before a fresh download.\\
\emph{What to do:} Always read the \prompt{Using ... file:} line. Keep one
protein per project folder; start a new project to switch proteins. See
Chapter~\ref{ch:projectfolder}.

{\bfseries \prompt{PDB XQ99 not found (404). May be a large structure only
available in mmCIF format.}}\\*
\emph{Why:} The code does not exist, or the entry is so large that the Protein
Data Bank offers it only in mmCIF format.\\
\emph{What to do:} Check the code on the RCSB website. Very large entries
cannot be loaded in version 1.23.0. See Chapter~\ref{ch:structure}.

{\bfseries \prompt{No AlphaFold structure found for P69905} for a protein that
has an AlphaFold model.}\\*
\emph{Why:} ProPrep reports a failed internet connection this way.\\
\emph{What to do:} Check your internet connection and try again. See
Chapter~\ref{ch:structure}.

{\bfseries An AlphaFold or AlphaFill file is not listed in the file browser,
or ProPrep says \prompt{AF-P69905-F1-model\_v6.cif is not a PDB file.}}\\*
\emph{Why:} ProPrep's file browsers accept only PDB-format files
(\file{.pdb}), not mmCIF files (\file{.cif}).\\
\emph{What to do:} Load the model directly by its UniProt accession, or
download the PDB-format file from the AlphaFold Database website. See
Chapters~\ref{ch:structure} and~\ref{ch:compare}.

{\bfseries \prompt{Enter a single number, '..', or 'q'.} when you type a file
name in a file browser.}\\*
\emph{Why:} The browser needs a number from its list, \texttt{..} to go up,
or a full path.\\
\emph{What to do:} Type the number of the file, or its full path. See
Chapter~\ref{ch:structure}.

{\bfseries ProPrep quits after \key{m} (manual matrix) in the Biological
Assembly Generator, with \prompt{Error: name 'self' is not defined}.}\\*
\emph{Why:} In version 1.23.0 this option does not work.\\
\emph{What to do:} Do not use \key{m}. Choose an assembly number, or download
the assembly directly in the Structure Loader (structure type~2). Resume your
session to get back to where you were. See Chapter~\ref{ch:structure}.

{\bfseries \prompt{No REMARK 350 records found in PDB file.} in \key{1b}.}\\*
\emph{Why:} \key{1b} is reading a file ProPrep made, such as the
hydrogen-stripped copy, which has no header with assembly information.\\
\emph{What to do:} Run \key{1b} in a new project right after loading the
structure, or download the assembly directly. See
Chapter~\ref{ch:structure}.

\section{Cleaning and repairing the structure}\label{sec:troubleshooting-cleaning}

{\bfseries \prompt{Error executing PDB Filter.filter\_pdb\_structure: invalid
literal for int() with base 10: '1-50'} (or \prompt{list index out of
range}); all your filter answers are lost.}\\*
\emph{Why:} At \prompt{Select residues} the PDB Filter (\key{4a}) does not
accept ranges such as \typed{1-50}, or a number that is not in the table.\\
\emph{What to do:} Run the filter again and type single numbers from the
table, separated by commas, or \typed{all}. See Chapter~\ref{ch:cleaning}.

{\bfseries \prompt{Received '2,4', which is not one of the available options}
at the water question of the PDB Filter.}\\*
\emph{Why:} Although the prompt mentions a list, only one choice is accepted
there.\\
\emph{What to do:} Type a single number. See Chapter~\ref{ch:cleaning}.

{\bfseries The Structure Fixer (\key{4c}) shows \prompt{MODELLER off: capping
only}; loops, missing side chains or mutations are not added.}\\*
\emph{Why:} Building missing parts, standard mutations and the choice between
alternate locations all need MODELLER with a valid license key.\\
\emph{What to do:} Add a MODELLER key (Chapter~\ref{ch:installation}).
Without one, tLEaP still adds missing side-chain atoms when it builds the
topology. See Chapter~\ref{ch:cleaning}.

{\bfseries \prompt{No terminal caps selected.} and nothing is written; a
pending mutation disappears.}\\*
\emph{Why:} Without MODELLER, the Structure Fixer writes a file (and applies a
non-standard mutation) only when at least one end cap is added.\\
\emph{What to do:} Add at least one cap (ACE or NME), or get a MODELLER key.
See Chapter~\ref{ch:cleaning}.

{\bfseries Residue numbers change after repair, for example 2RH1 residues
29--1161 become 1--506.}\\*
\emph{Why:} The Structure Fixer numbers residues from 1 when it writes the
repaired file. An added ACE cap shifts the numbers by one more.\\
\emph{What to do:} Use the new numbers in every later step. See
Chapters~\ref{ch:cleaning} and~\ref{ch:protonation}.

{\bfseries The template comparison in \key{4c} reports a missing OXT atom on
almost every residue.}\\*
\emph{Why:} These are false alarms: OXT exists only on the last residue of a
chain.\\
\emph{What to do:} Use the list of missing atoms from REMARK~470 instead. See
Chapter~\ref{ch:cleaning}.

\section{Protonation and orientation}\label{sec:troubleshooting-protonation}

{\bfseries \prompt{Invalid pH value, using default of 7.0} (or \prompt{Invalid
threshold, using default of 0.5}) and ProPrep carries on.}\\*
\emph{Why:} The Protonation State Analyzer (\key{4e}) does not ask again
after a typing mistake.\\
\emph{What to do:} Run option \key{1} again with the value you want. See
Chapter~\ref{ch:protonation}.

{\bfseries \prompt{Skipping invalid syntax: XYZ (expected resid:NAME)} or
\prompt{Invalid name 'LYS' - valid choices: GLH, GLU}, and ProPrep moves
on.}\\*
\emph{Why:} An invalid override is skipped without a second chance; the
recommendation is used instead.\\
\emph{What to do:} Run option \key{3} again and type the change in the form
shown, for example \typed{21:GLH}. See Chapter~\ref{ch:protonation}.

{\bfseries The Structure Orientator (\key{4f}) keeps using an old file after
you filtered or protonated again; file names grow to
\file{..._oriented_oriented_...pdb}.}\\*
\emph{Why:} In \key{4f}, an oriented file comes before every other file, so
each run starts from the previous result.\\
\emph{What to do:} Orient as the last preparation step. If you have to change
the structure afterwards, start a new project. See
Chapter~\ref{ch:protonation}.

\section{Redox sites, force fields and QM/MM}\label{sec:troubleshooting-redox}

{\bfseries ACE, GOL, ACT or EDO are listed as \prompt{organic cofactor} by the
Redox Site Detector (\key{2a}).}\\*
\emph{Why:} Every non-standard residue is treated as a possible cofactor,
including caps and additives from crystallization.\\
\emph{What to do:} At \prompt{Select additional centers to include}, type
only the numbers you want. See Chapter~\ref{ch:redox}.

{\bfseries A bond of 8.21\,Å was created without a warning.}\\*
\emph{Why:} The list of nearby atoms is sorted again for every source atom, so
a number from an earlier list can point to a different atom.\\
\emph{What to do:} Type atom names such as \typed{SG} instead of numbers. To
repair, type the source atom again and choose \textsf{2 Replace one existing
bond}. See Chapter~\ref{ch:redox}.

{\bfseries The Redox Site Preparer (\key{4d}) shows \texttt{None} under
\textsf{Compatible Transformers} for a single heme site.}\\*
\emph{Why:} Some bonds that the heme recipe needs inside the heme were not
recorded, because the single-site route in \key{2a} does not show the bond
requirements.\\
\emph{What to do:} In \key{4d} type \typed{why 1} to see what is missing, then
add those bonds in \key{2a}. See Chapter~\ref{ch:redox}.

{\bfseries The site editor in \key{2a} drops back to the menu with
\prompt{Error handling menu option 'edit': invalid literal for int() with base
10}; your edits are lost.}\\*
\emph{Why:} Any text that is not a number at a residue-number prompt ends the
editor.\\
\emph{What to do:} Give one answer per prompt and type \typed{save} often. See
Chapter~\ref{ch:redox}.

{\bfseries All options of \key{4d} are greyed out with \prompt{Already
generated microstate(s)}.}\\*
\emph{Why:} Microstates can be generated only once per session.\\
\emph{What to do:} Start a new session in a new project folder. See
Chapter~\ref{ch:redox}.

{\bfseries \prompt{Required center count (expected 1, got 4)} for
hemoglobin (4HHB).}\\*
\emph{Why:} Several hemes were grouped into one site, but each heme recipe
expects one heme per site.\\
\emph{What to do:} Detect the hemes again as separate sites, one heme each.
See Chapter~\ref{ch:redox}.

{\bfseries \prompt{'FD1' is not a valid option for site 1} for an
iron--sulfur cluster.}\\*
\emph{Why:} The broken-symmetry code must belong to the redox state chosen
for that site (FO, FR or FD codes).\\
\emph{What to do:} Use a code from the list ProPrep prints with the message.
See Chapter~\ref{ch:redox}.

{\bfseries \prompt{Step 4 failed: Gaussian output not found: mk1.log}
followed by a long technical report, in the Force Field Parameterizer
(\key{2c}).}\\*
\emph{Why:} This step needs the result of a Gaussian calculation that you run
yourself. ProPrep has not crashed.\\
\emph{What to do:} Type \key{s} to save and quit, run Gaussian on the
\file{.gjf} file, then resume with \key{2c} option \key{3}. Without Gaussian,
make AM1-BCC charges outside ProPrep and import them. See
Chapter~\ref{ch:forcefields}.

{\bfseries \prompt{Frame extraction failed: No prmtop file selected} in the
QM/MM Preparator (\key{5c}).}\\*
\emph{Why:} \key{5c} starts from a topology, and none was found.\\
\emph{What to do:} Build the topology with \key{5a} first. See
Chapter~\ref{ch:qmmm}.

{\bfseries \prompt{orca\_mm not found ... Please ensure ORCA is installed}
and no files are written (ORCA route of \key{5c}).}\\*
\emph{Why:} ORCA's helper program \texttt{orca\_mm} is needed even to write the
input files.\\
\emph{What to do:} Install ORCA and give ProPrep its folder, or use the
Gaussian route. See Chapter~\ref{ch:qmmm} and Appendix~\ref{app:software}.

\section{Building the topology (tLEaP)}\label{sec:troubleshooting-topology}

{\bfseries Back at the Topology Generator menu after \prompt{Please run
RedoxSite detection and try again.}; no \file{_tleap.in} file was written.}\\*
\emph{Why:} You pressed \key{Enter} at \prompt{Do you want to continue without
bond definitions? [y/n] (n)}. The default \key{n} stops the build.\\
\emph{What to do:} Run \key{5a} option \key{2} again (your settings are
remembered) and type \key{y}. For an ordinary protein no bond definitions are
needed. See Chapter~\ref{ch:topology}.

{\bfseries \prompt{tLEaP completed successfully}, then \prompt{Errors = 42};
or the menu says the topology was generated, but the \file{.prmtop} file is
empty.}\\*
\emph{Why:} In version 1.23.0 the success messages are printed even when
tLEaP fails, and a 0-byte topology counts as present.\\
\emph{What to do:} Trust \prompt{Errors =} and the file size. Read
\file{leap.log} for the cause. See Chapter~\ref{ch:topology}.

{\bfseries \prompt{tLEaP failed (exit code: 31)} with only \prompt{Unknown
residue: HOH} on the screen.}\\*
\emph{Why:} Crystal waters are left in an implicit-solvent build. The real
reason, \prompt{Atom ... does not have a type}, is only in \file{leap.log}.\\
\emph{What to do:} Remove the waters with the PDB Filter (\key{4a}) before
building without explicit water. See Chapter~\ref{ch:topology}.

{\bfseries \prompt{Unknown residue: HEC} (or another cofactor name), tLEaP
fails with many errors.}\\*
\emph{Why:} The heme, ligand or metal has no force-field parameters yet.
ProPrep does not warn you before the build.\\
\emph{What to do:} Prepare it first: hemes and other redox cofactors with
\key{2a} and \key{4d} (Chapter~\ref{ch:redox}), other molecules with \key{2c}
(Chapter~\ref{ch:forcefields}).

{\bfseries \prompt{Could not find bond parameter for atom types: SH -
SH}.}\\*
\emph{Why:} Cysteines in disulfide bonds are still named CYS. The
Protonation State Analyzer does not rename them to CYX.\\
\emph{What to do:} Rename them to CYX (\key{2a} then \key{4d}) \emph{and} add
a \texttt{bond} line for each pair in \key{5a} option \key{1}. See
Chapter~\ref{ch:topology}.

{\bfseries \prompt{Detected CYX residue with a Sulfur atom not bonded to
another CYX Sulfur! Did you forget to make the disulfide bond?}}\\*
\emph{Why:} The cysteines are named CYX, but no bond lines were given, so the
topology has no disulfide bonds although the build ``succeeded''.\\
\emph{What to do:} Add the bond lines (\key{5a} option \key{1}), then
regenerate the input (\key{2}) and the topology (\key{4}). See
Chapter~\ref{ch:topology}.

{\bfseries A good \file{.prmtop} file became 0 bytes after you built a second
version (for example implicit solvent).}\\*
\emph{Why:} Both tLEaP inputs write the same output names. A failed second
build empties the first file, without a warning.\\
\emph{What to do:} Give each version its own output names on the
\texttt{saveamberparm} line when you edit the template. See
Chapter~\ref{ch:topology}.

{\bfseries The constant-pH input file (\file{.cpin}) lists only LYS and TYR;
ASP, GLU and HIS are missing.}\\*
\emph{Why:} The topology was not built with the constant-pH force field, or
the residues were not renamed (AS4, GL4). There is no warning.\\
\emph{What to do:} Rename the residues with the constant-pH option of
\key{4e} and build with the \textsf{Constant pH (ff10)} protein force field.
See Chapter~\ref{ch:topology}.

\section{Running a simulation}\label{sec:troubleshooting-md}

{\bfseries \prompt{Error starting simulation: [Errno 2] No such file or
directory: 'pmemd'} and \prompt{Simulations started: 0}.}\\*
\emph{Why:} The default engine is \texttt{pmemd}, which is not installed with
ProPrep. The recommended assignment even picks \texttt{pmemd.cuda} on a
computer without a GPU.\\
\emph{What to do:} Go back to \key{5b} option \key{1}. At the hardware step,
choose \key{3} (same engine for all) and then \texttt{sander} (or
\texttt{sander.MPI}). The queue menu of the Execute screen cannot change the
engine. See Chapter~\ref{ch:runningmd}.

{\bfseries The review table still shows the original lengths after you
shortened the stages.}\\*
\emph{Why:} Only the display is not updated.\\
\emph{What to do:} Nothing; your changes are in the files. Open
\file{simulation.mdin} in a stage folder to confirm. See
Chapter~\ref{ch:runningmd}.

{\bfseries \prompt{Enable Constant pH MD?} is offered with the default
\key{y}, and \key{5b} uses a \file{_cpin.prmtop} topology.}\\*
\emph{Why:} Making a cpin file in \key{5a} switches the default structure pair
to the constant-pH one.\\
\emph{What to do:} For an ordinary simulation, answer \key{n}, and select the
normal topology and coordinate pair. See Chapter~\ref{ch:runningmd}.

{\bfseries \prompt{Error importing restraints: 'MDRestraintModule' object has
no attribute 'get\_from\_workspace'}.}\\*
\emph{Why:} The message appears although the import worked.\\
\emph{What to do:} Check with option \key{3} (display); if your restraints are
listed, carry on. See Chapter~\ref{ch:runningmd}.

\section{Analyzing results}\label{sec:troubleshooting-analysis}

{\bfseries \prompt{Returning to MD Manager menu...} right after you type
your list of analyses.}\\*
\emph{Why:} In version 1.23.0 any list that contains the digit 0, such as
\typed{10}, \typed{20} or \typed{15,18,20}, is read as ``exit''.\\
\emph{What to do:} Leave out analyses 10 (angle) and 20 (density maps); they
cannot be run in this version. See Chapter~\ref{ch:analysis}.

{\bfseries \prompt{Failed to setup action} and a long technical report; you
are back at the MD Manager menu.}\\*
\emph{Why:} An atom selection does not exist, for example the CB atom of a
glycine in a dihedral analysis.\\
\emph{What to do:} Start the analysis again and choose atoms that exist. See
Chapter~\ref{ch:analysis}.

\section{Docking}\label{sec:troubleshooting-docking}

{\bfseries \prompt{Template matching failed for: ['A:99', 'B:99']} or
\prompt{Creation of data structure for receptor failed}.}\\*
\emph{Why:} The last residue of each chain has no OXT atom, but ProPrep set
it up as a charged chain end. This is common in PDB files.\\
\emph{What to do:} In option \key{4}, type the residue (for example
\typed{A:99}) and answer \key{b} (chain break); repeat for each chain. See
Chapter~\ref{ch:docking}.

{\bfseries Docking finishes in under a second with weak scores (about
$-4$\,kcal/mol).}\\*
\emph{Why:} The search box was placed around a ligand built from SMILES,
which sits at the origin, not in the protein.\\
\emph{What to do:} Define the box around named residues (\key{r}) or by
typing its center and size (\key{c}). See Chapter~\ref{ch:docking}.

{\bfseries \prompt{The ligand could not be built: Unsupported ligand file type
'.smi'; use .sdf, .mol or .mol2.}}\\*
\emph{Why:} A \file{.smi} library works for a campaign, but not as a single
ligand. The dashboard still shows its name.\\
\emph{What to do:} Run a campaign, or use an \file{.sdf}, \file{.mol} or
\file{.mol2} file. Set the search box again after every ligand change. See
Chapter~\ref{ch:docking}.

\section{Membrane systems}\label{sec:troubleshooting-membrane}

{\bfseries \prompt{tLEaP could not add the hydrogens ... Build anyway with
packmol-memgen's own protonation?}; if you answer \key{y}, the build is
refused several minutes later with \prompt{The protein in bilayer.pdb is not
the one packmol-memgen was given}.}\\*
\emph{Why:} The structure still contains ions, crystal waters or other
residues that tLEaP cannot handle.\\
\emph{What to do:} Answer \key{n}. Remove every non-protein component with the
PDB Filter (\key{4a}), then build again. See Chapter~\ref{ch:membrane}.

{\bfseries Lipid, box or orientation settings are back at their defaults when
you open \key{4g} again.}\\*
\emph{Why:} Choosing another protein, filtering, or returning after a build
clears them without a message.\\
\emph{What to do:} Read the \textsf{Current Configuration} block, or press
\key{r} (review), before every build. See Chapter~\ref{ch:membrane}.

{\bfseries Hundreds of tLEaP warnings scroll past after a membrane build.}\\*
\emph{Why:} Lipid renaming and close contacts from packing produce many
warnings.\\
\emph{What to do:} Nothing, as long as the summary says \prompt{Errors = 0}.
The first minimization removes the close contacts. See
Chapter~\ref{ch:membrane}.

\section{Sessions, batch runs and the browser}\label{sec:troubleshooting-sessions}

{\bfseries ProPrep replays your whole session and then quits before you can
type anything.}\\*
\emph{Why:} You chose \textsf{Continue from end} on a session that ended with
a normal exit, so the exit is replayed too.\\
\emph{What to do:} Open the Session Editor and delete the last two answers,
for example \typed{d 18-end}, then save. Use \textsf{Continue from end} only for
interrupted sessions. See Chapter~\ref{ch:sessions}.

{\bfseries \prompt{Invalid choice: undo. Please try again.}}\\*
\emph{Why:} Session recording is off (\texttt{--no-session}, or option~3 at
start-up).\\
\emph{What to do:} Start ProPrep again with recording on. See
Chapter~\ref{ch:sessions}.

{\bfseries After \typed{undo} and \key{c} (change), ProPrep keeps the old
answer: \prompt{[REPLAY: 'r' now means a different option; using 'd' ...]}.}\\*
\emph{Why:} For menu questions, the stored answer is not updated in version
1.23.0.\\
\emph{What to do:} Use \key{r} (re-answer live) after \typed{undo}, or change
the answer in the Session Editor. See Chapter~\ref{ch:sessions}.

{\bfseries A replay stops to ask \prompt{Select menu mode [1/2]}.}\\*
\emph{Why:} That question comes before recording starts, so it is not in the
session file.\\
\emph{What to do:} Add \texttt{--full-menu} (or \texttt{--workflow}) to the
replay command. See Chapters~\ref{ch:sessions} and~\ref{ch:cmdline}.

{\bfseries \prompt{The recorded session was started with --pdbfile
my\_protein.pdb, but that file is not here}.}\\*
\emph{Why:} The structure file the recording started from has been moved or
deleted. Without it every tool says \prompt{Needs a loaded structure}.\\
\emph{What to do:} Copy the file back into the project folder, or replay with
\texttt{--pdbfile} and the file's path. See Chapter~\ref{ch:sessions}.

{\bfseries Every batch run fails with \prompt{Setup error: [Errno 2] No such
file or directory: 'batch\_out/\allowbreak run\_01'}.}\\*
\emph{Why:} \texttt{--output-dir} does not create the output folder.\\
\emph{What to do:} Create the folder first (\texttt{mkdir batch\_out}) and
run the batch again. See Chapter~\ref{ch:batch}.

{\bfseries \prompt{Replay diverged: no recorded answer for prompt ...} for one
protein of a batch.}\\*
\emph{Why:} That protein is asked different questions than the recorded one.
For example, 1CRN has no crystal waters, so the water question never comes.\\
\emph{What to do:} Use proteins with the same make-up (chains, waters,
ligands), or record a template that fits them. Add \texttt{--batch-continue}
to finish the others. See Chapter~\ref{ch:batch}.

{\bfseries \prompt{! Input column 'P69905' matches no template variable} for a
plain list of UniProt accessions.}\\*
\emph{Why:} A first line that starts with a letter is taken as a column
heading.\\
\emph{What to do:} Put the line \typed{input\_protein} at the top of the list
file. See Chapter~\ref{ch:batch}.

\section{General: an error message full of technical text}\label{sec:troubleshooting-general}

{\bfseries Many lines of technical text appear, often starting with
\prompt{Traceback (most recent call last)} or \prompt{Error executing ...},
and ending with a line such as \prompt{ValueError: ...}.}\\*
\emph{Why:} This is a \term{traceback}: the program's own report of where an
error happened. It is meant for the developers. In most cases ProPrep has not
crashed: it stops the current action and returns to a menu.\\
\emph{What to do:}
\begin{enumerate}
  \item Read the \emph{last} line. It often names the cause in plain words,
  such as a file that is missing or a value it could not read.
  \item Look up that message in this chapter.
  \item Check what was saved. Answers given within the failed action are
  usually lost; files written before it are kept.
  \item Repeat the step, carefully, with valid answers. If it fails again,
  start ProPrep with \texttt{--debug} and report the problem with the
  error lines and your session file (Chapter~\ref{ch:help}).
\end{enumerate}

{\bfseries Lines with a date, a time and \texttt{INFO}, \texttt{WARNING} or
\texttt{ERROR} appear in the middle of the screen.}\\*
\emph{Why:} Some of ProPrep's internal log messages are shown on the screen.\\
\emph{What to do:} Usually nothing. Read on to ProPrep's own message below
them. See Chapter~\ref{ch:analysis}.

{\bfseries A success message, but the expected file is missing or
empty.}\\*
\emph{Why:} Some steps in version 1.23.0 report success without checking the
result (Chapter~\ref{ch:knownissues}).\\
\emph{What to do:} Look in the project folder and check the file size, for
example with \texttt{ls -l} in the Terminal. A 0-byte file is a failed step.
See Chapter~\ref{ch:projectfolder}.

{\bfseries You pressed \key{Enter} in the feedback form and it tried to post
your report.}\\*
\emph{Why:} The default of \prompt{Submit this feedback to GitHub Issues?
[y/n] (y)} is yes.\\
\emph{What to do:} Type \key{n} at that question if you are not ready, or
press \key{Enter} at \prompt{Summary:} to cancel the form. See
Chapter~\ref{ch:help}.
